% Resposta da questão 15

Considere que a velocidade do som $c$ dependa da pressão $p$ e da densidade $\rho$ segundo uma lei do tipo

\begin{equation}
    c \propto p^{a}\rho^{b}.    
\end{equation}

As dimensões das grandezas são:

\begin{equation}
   [c]=LT^{-1}, 
\qquad 
    [p]=ML^{-1}T^{-2}, 
\qquad 
    [\rho]=ML^{-3}. 
\end{equation}

Substituindo na relação dimensional:

\begin{equation}
     LT^{-1}=(ML^{-1}T^{-2})^{a}(ML^{-3})^{b}.  
\end{equation}

Expandindo:

\begin{equation}
    LT^{-1}=M^{a+b}L^{-a-3b}T^{-2a}.   
\end{equation}


Igualando os expoentes das dimensões fundamentais:

Para $M$:
\begin{equation}
        a + b = 0
\end{equation}

Para $L$:
\begin{equation}
        -a - 3b = 1
\end{equation}

Para $T$:
\begin{equation}
        -2a = -1
\end{equation}

Da última equação:
\begin{equation}
        a = \frac{1}{2}
\end{equation}

Substituindo em $a + b = 0$:
\begin{equation}
        b = -\frac{1}{2}
\end{equation}

Logo:

\begin{equation}
    c \propto p^{1/2}\rho^{-1/2}.
\end{equation}

Portanto, a relação entre as grandezas é:

\begin{equation}
    c = k\sqrt{\frac{p}{\rho}},  
\end{equation}

\noindent onde \(k\) é uma constante adimensional.

Assim, a velocidade do som é proporcional à raiz quadrada da razão entre a pressão e a densidade do gás.